\documentclass{report}
\usepackage{amsmath,amssymb}
\setlength{\parindent}{0mm}
\setlength{\parskip}{1em}
\begin{document}
\begin{center}
\rule{6in}{1pt} \
{\large Sara Pollock \\
{\bf Regularization and Adaptivity for Nonlinear Elliptic PDE Solvers}}

Department of Mathematics \\ 3368 TAMU \\ College Station \\ TX 77843-3368
\\
{\tt snpolloc@math.tamu.edu}\end{center}

In this talk I will introduce an adaptive framework developed to solve
nonlinear elliptic partial differential equations (PDE) starting from a
coarse mesh. The target problem class includes quasi-linear problems with
steep internal layers in the solution-dependent diffusion coefficients,
for which standard methods such as Newton or Picard iterations are known
to fail. The method is designed to start with a discretization that does
not resolve the problem coefficients. The discrete problem on the initial
sequence of meshes is not assumed to inherit the stability, coercivity,
monotonicity or solvability properties of the continuous system;
essentially, the initial sequence of discrete problems is assumed
ill-posed.

A sequence of partial solves of regularized problems is used to refine
the discretization where necessary to resolve the problem coefficients
and data. Adaptivity is used both for mesh refinement and automatic
control of the regularization parameters to ultimately solve the discrete
problem without regularization. I will discuss improving stability of the
method with different pseudo-time integrators. The method will be
demonstrated with numerical examples using an underlying linear finite
element discretization.


\end{document}
